\chapter{Requirement Specifications}\label{ch:req}

% Local helper for the use-case description tables (used only in this chapter).
% #1 label suffix, #2 ID, #3 name, #4 actors, #5 precondition, #6 postcondition,
% #7 main flow, #8 alternative flow, #9 exception flow
\newcommand{\rclUseCase}[9]{%
\begin{table}[!htbp]
\centering
\small
\caption{Textual description of use case #2: #3}\label{tab:uc#1}\label{uc:#1}
\begin{tabularx}{\textwidth}{@{}>{\bfseries\raggedright\arraybackslash}p{2.7cm} X@{}}
\toprule
Use Case ID & #2\\
Name & #3\\
Actor(s) & #4\\
Precondition & #5\\
Postcondition & #6\\
Main flow & #7\\
Alternative flow & #8\\
Exception flow & #9\\
\bottomrule
\end{tabularx}
\end{table}}

This chapter specifies what RecaLL must do and under which conditions. The requirements
were traced from the delivered code: each use case below names the routes that implement
it, taken from the route inventory of the backend (158 route handlers: 87 under
\texttt{/api/v1} in \texttt{app/api/v1/*.py} and 71 older ``study'' routes in
\texttt{app/main.py}).

\section{Stakeholders for the System}

\Cref{tab:stakeholders} lists the people and organisations with an interest in the system.

\begin{small}
\begin{longtable}{@{}>{\raggedright\arraybackslash}p{3.0cm} >{\raggedright\arraybackslash}p{5.4cm} >{\raggedright\arraybackslash}p{5.6cm}@{}}
\caption{Stakeholders}\label{tab:stakeholders}\\
\toprule
\textbf{Stakeholder} & \textbf{Role} & \textbf{Interest in the system / influence}\\
\midrule
\endfirsthead
\toprule
\textbf{Stakeholder} & \textbf{Role} & \textbf{Interest in the system / influence}\\
\midrule
\endhead
\bottomrule
\endfoot
Student (primary user) & Imports course material, reads, makes notes, takes drills and
quizzes, reviews flashcards, chats with the Knowledge Hub & Better retention for less effort;
privacy of their material; low or zero cost. High influence: the only end-user role\\
Self-learner / professional & Same as the student, outside a university course & Studying
technical material on their own computer\\
Course teachers (customer) & Set the project requirements and evaluate the
product and report & A complete software-engineering process and a working, tested product.
High influence\\
Development team & Builds, tests, packages and maintains RecaLL & Maintainable code, clear
requirements, manageable scope\\
Cloud AI providers (Google Gemini, OpenAI, Groq, OpenRouter) & Run the language models that
RecaLL calls through the student's key & Terms of use, quotas and rate limits. Low direct
influence, but outages and quota changes affect every AI feature\\
Local model ecosystem (Ollama and open-weight models) & Runs models on the student's machine &
Hardware limits and model quality; enables fully offline use\\
Optional services (Supabase, Langfuse) & Cloud sync and sign-in; LLM tracing & Only involved
when configured; off in the delivered build\\
Authors and publishers of study material & Own the copyright of uploaded textbooks and
slides & Material must not be redistributed; RecaLL keeps it on the student's computer\\
CUET and other institutions & Host the course and could recommend the tool & Academic
integrity and responsible use of AI\\
Regulators & Bangladeshi data-protection authority (PDPO 2025 \cite{pdpo2025}); GDPR
authorities for EU users \cite{gdpr2016} & Lawful processing of personal data\\
\end{longtable}
\end{small}

\section{Use Case Diagram with Graphical and Textual Description}

\Cref{fig:usecase} shows the use cases. The primary actor is the \emph{Student}. The
\emph{AI provider} (a cloud model or a local Ollama server) is a secondary actor in every
use case that generates or grades content, and \emph{Supabase} is an optional secondary
actor for sign-in and sync. There is no administrator role: RecaLL is a single-user desktop
application, and each account owns exactly one personal workspace.

\widefig{usecase.pdf}{Use case diagram of RecaLL}{fig:usecase}

The textual descriptions of all twelve use cases follow in
\Cref{tab:uc01,tab:uc02,tab:uc03,tab:uc04,tab:uc05,tab:uc06,tab:uc07,tab:uc08,tab:uc09,tab:uc10,tab:uc11,tab:uc12}.
HTTP status codes in the exception flows are the ones the backend returns
(\texttt{app/core/errors.py}).

\rclUseCase{01}{UC-01}{Register / log in}
{Primary: Student. Secondary: Supabase (optional)}
{The application is running and the student is not signed in.}
{The student holds a signed JWT (HS256, valid for 60 minutes) linked to a personal workspace,
and the backend has switched the study database to \texttt{user\_<id>\_recall.db}.}
{1.~The student opens \texttt{/register} and enters name, e-mail and password.
2.~The interface calls \texttt{POST /api/v1/auth/signup}. 3.~The system checks the rate
limit (15 requests per minute per address), validates the e-mail format and password length
(8--128 characters), hashes the password with PBKDF2-HMAC-SHA256 and, in one transaction,
creates the user, a personal workspace and default preferences. 4.~The system returns a JWT
carrying the user and workspace IDs. 5.~The interface calls \texttt{POST /api/set-user}, which
selects the student's own study database. 6.~The dashboard opens.}
{1a.~A returning student opens \texttt{/login}; \texttt{POST /api/v1/auth/login} verifies the
hash in constant time and returns a JWT. 1b.~If Supabase is configured, the student signs in
through Supabase and its token is accepted by the backend. 1c.~A signed-in student changes the
password through \texttt{POST /api/v1/auth/change-password}.}
{E1.~E-mail already registered: 409. E2.~Wrong e-mail or password: 401 with one uniform
message (no account enumeration). E3.~Account deactivated: 401. E4.~Too many attempts: 429.
E5.~Invalid input: 422 with field details.}

\rclUseCase{02}{UC-02}{Upload and process a document}
{Primary: Student. Secondary: AI provider (topic processing), Marker (optional OCR)}
{The student is signed in. The file is a PDF, Markdown or text file of at most 50~MB.}
{The document is stored, its topic tree is saved in the study database, and the document and
its chunks are available in the application database for chat, flashcards and quizzes.}
{1.~The student chooses \emph{Add Knowledge} and selects a file. 2.~The interface computes a
SHA-256 hash and sends \texttt{POST /books/upload}. 3.~The system checks the size, recomputes
the hash on the server and stores the PDF under its hash. 4.~The system reuses an existing book
with the same hash for this user, or creates a new one. 5.~The system reads the PDF's built-in
table of contents, falls back to heading detection, filters front-matter entries and saves the
topic tree. 6.~The system mirrors the book into \texttt{recall\_saas.db}: it chunks the text
(1{,}500 characters with 150 overlap), embeds the chunks and saves them. 7.~The student picks
topics, and \texttt{POST /books/\{id\}/process-stream} processes them, streaming progress as
SSE events.}
{2a.~Through the API, \texttt{POST /api/v1/documents/upload} validates the file, creates a
document and a processing job, and runs the six-stage pipeline (validation, extraction,
normalisation, chunking, embedding, indexing) as a background task; this pipeline extracts PDF,
Markdown and text directly with PyMuPDF. The interface polls
\texttt{/documents/\{id\}/status}. 5a.~On the book path only, a scanned PDF of up to 25 pages is
sent to the Marker extractor when it is installed or a Marker (Datalab) API key is set; larger
scanned files can be re-parsed later (\texttt{/books/\{id\}/reparse-handwriting-stream}).}
{E1.~File larger than the limit: 413. E2.~Unsupported format or missing \texttt{\%PDF}
signature: 422. E3.~Document or storage quota of the plan reached: 402. E4.~Extraction fails:
the job and document are marked \texttt{failed} with the error text. E5.~Scanned PDF and no
Marker: no topics are produced, and the student is told to configure OCR.}

\rclUseCase{03}{UC-03}{Read and annotate a document}
{Primary: Student. Secondary: AI provider (explain, flashcards)}
{A document has been imported (UC-02).}
{Highlights and notes are saved per book and page; optional explanations or flashcards are
created from a selected passage.}
{1.~The student opens the document; the reader shows the topic tree and the original PDF pages
(\texttt{GET /books/\{id\}/pdf}). 2.~The student selects a passage and creates a highlight or
note (\texttt{POST /books/\{id\}/annotations}). 3.~The system stores the annotation with its page
and position. 4.~The reader shows all annotations of the page
(\texttt{GET /books/\{id\}/annotations?page=n}). 5.~The student saves the reading position
(\texttt{PATCH /books/\{id\}/reading-state}).}
{2a.~The student asks for an explanation of the passage
(\texttt{POST /books/\{id\}/annotations/explain}). 2b.~The student generates flashcards from the
passage (\texttt{/annotations/generate-flashcards}). 4a.~The student edits or deletes an
annotation. 5a.~The student exports the annotated PDF (\texttt{/books/\{id\}/export-annotated}).}
{E1.~Book or PDF file missing: 404. E2.~AI provider fails during an explanation: an error
message is shown and the annotation itself is kept.}

\rclUseCase{04}{UC-04}{Generate Cornell notes}
{Primary: Student. Secondary: AI provider}
{A document with at least one topic exists; an AI provider is configured.}
{A Cornell note with seven fixed parts is stored for the topic and shown in the notebook
editor, where the student can edit it.}
{1.~The student opens a topic and presses \emph{Generate Cornell Note}. 2.~The interface calls
\texttt{POST /topics/\{id\}/notes/scaffold-stream}. 3.~The system loads the topic text (waiting
for any running concept extraction to finish) and fills the \texttt{cornell\_scaffold\_prompt}
template. 4.~The provider
generates the note: core intuition, invariants and definitions, step-by-step mechanisms, a
worked example, visual comparisons, exam pitfalls and self-test cue questions. 5.~Progress and
the result stream to the interface as SSE events. 6.~The note is saved against the topic and
rendered with KaTeX.}
{4a.~For a parent topic, a note is generated for each leaf subtopic (in parallel, within the
provider's concurrency limit, reusing notes already saved) and the results are combined into one
master guide.
6a.~The student edits the note, adds images (\texttt{POST /notes/upload-image}) or appends
text (\texttt{/topics/\{id\}/notes/append}). 6b.~The student deletes the note.}
{E1.~Topic not found: 404. E2.~Provider error, timeout or rate limit: an SSE error event is
sent and no note is saved.}

\rclUseCase{05}{UC-05}{Take a Socratic drill}
{Primary: Student. Secondary: AI provider}
{A topic (optionally one extracted concept) with processed text exists; an AI provider is
configured.}
{The attempt and its score are recorded, the topic's (and concept's) mastery score is updated,
and any cards the student accepted are added to the review deck.}
{1.~The student starts a drill on a topic or concept. 2.~\texttt{POST /topics/\{id\}/drill/generate}
fills \texttt{socratic\_question\_prompt}; the provider returns exactly two questions (the first a
causal mechanism, the second a counterfactual or applied scenario), each with key invariants
and a hint. 3.~The reply is repaired if needed and validated against a schema. 4.~The student
writes an explanation. 5.~\texttt{POST /topics/\{id\}/drill/evaluate} fills
\texttt{socratic\_evaluation\_prompt} with the question, invariants, topic text and answer.
6.~The system returns a mastery score (0--100) and band, knowledge gaps, misconceptions, a
Socratic nudge and 0, 1 or 1--2 suggested flashcards. 7.~The student selects cards to keep,
and \texttt{POST /topics/\{id\}/drill/save-cards} saves them.}
{6a.~Score 85 or above: no cards are suggested. 6b.~Score below 60: no nudge is given, because
the student needs to re-study first. 7a.~The student pins the gaps to the topic notes instead
of saving cards.}
{E1.~Topic not found: 404. E2.~The provider's JSON cannot be repaired or validated: an error is
returned to the widget and the student can retry.}

\rclUseCase{06}{UC-06}{Review due flashcards (FSRS)}
{Primary: Student}
{At least one flashcard exists and is due.}
{Each rated card has a new FSRS state (difficulty, stability, retrievability) and due date, and
the review is logged for analytics.}
{1.~The student opens \emph{Review}. 2.~\texttt{POST /api/v1/reviews/sessions} starts a session over
the due queue. 3.~\texttt{GET /reviews/sessions/\{id\}/next} returns the next item; the question is
shown. 4.~The student reveals the answer (\texttt{POST /reviews/\{rid\}/reveal}). 5.~The student
rates the card Again, Hard, Good or Easy (\texttt{POST /reviews/\{rid\}/rate}). 6.~The
\texttt{FSRSScheduler} (desired retention 0.9, maximum interval 36{,}500 days) computes the next due
date. 7.~Steps 3--6 repeat until no item is left, and the session is completed.}
{7a.~The student abandons the session (\texttt{/abandon}). Note: the Review page uses only the
\texttt{/api/v1/reviews/*} routes above. The older study routes \texttt{GET /flashcards/due} and
\texttt{POST /flashcards/\{id\}/review} (ratings 1--4, with undo) also apply the \texttt{fsrs} library
to the cards in the study database, but the review interface does not call them.}
{E1.~Nothing is due: the queue is empty and the page says so. E2.~Session or review belongs to
another workspace: 404. E3.~Invalid rating: 422.}

\rclUseCase{07}{UC-07}{Ask the Knowledge Hub}
{Primary: Student. Secondary: AI provider}
{The student is signed in and has at least one processed document.}
{The question and the answer, with verified source citations, are saved in the conversation.}
{1.~The student opens the Knowledge Hub, picks or creates a conversation
(\texttt{POST /api/v1/conversations}) and optionally limits the scope to some documents.
2.~The student asks a question; \texttt{POST /conversations/\{id\}/messages} with
\texttt{stream=true}. 3.~The system checks that the conversation belongs to the workspace,
resolves the provider (saved BYOK key, then server key) and reserves usage credits (none for a
BYOK key). 4.~Hybrid keyword and vector search returns the five best chunks. 5.~The system builds
the prompt: system rules, up to six history messages and the chunks as
\texttt{<reference\_data id="S1">} blocks within 3{,}000 tokens. 6.~Tokens stream to the
interface as SSE \texttt{token} events. 7.~The citation validator keeps only \texttt{[S\emph{n}]}
references that point to retrieved chunks. 8.~The answer and sources are saved and a
\texttt{done} event is sent.}
{2a.~Without \texttt{stream}, the full answer is returned in one response. 8a.~From the answer, the
student uses \emph{Make Flashcard} or \emph{Test Knowledge} to move into UC-12 or UC-08.}
{E1.~Conversation not found or owned by another workspace: 404. E2.~Credits exhausted: 402.
E3.~Provider error: an SSE \texttt{error} event; the reserved usage is released.}

\rclUseCase{08}{UC-08}{Generate and take a quiz}
{Primary: Student. Secondary: AI provider}
{At least one processed document exists.}
{A quiz is stored; each attempt is stored with its answers, score and percentage.}
{1.~The student chooses documents, difficulty (easy, medium or hard), question types
(multiple choice, true/false) and a count of 1--50. 2.~\texttt{POST /api/v1/quizzes/generate}
retrieves chunks from the chosen documents and asks the provider for questions. 3.~The system
validates the questions and saves the quiz. 4.~The student starts an attempt
(\texttt{POST /quizzes/\{id\}/attempts}). 5.~Each answer is saved
(\texttt{POST /quiz-attempts/\{aid\}/answers}). 6.~On submit, the server scores all answers
deterministically. 7.~The results page shows the score and the correct answers.}
{1a.~The student edits or deletes questions, or deletes the quiz. 7a.~The student opens earlier
attempts (\texttt{GET /quizzes/\{id\}/attempts}).}
{E1.~No usable chunks in the chosen documents: 422. E2.~Quota exceeded: 402. E3.~Invalid
generated questions are discarded; if none remain, generation fails with an error.}

\rclUseCase{09}{UC-09}{View learning progress}
{Primary: Student}
{The student is signed in.}
{None (read-only).}
{1.~The student opens the Dashboard; \texttt{GET /api/v1/learning/dashboard} returns reviews due
and a knowledge summary. 2.~The student opens \emph{Progress}. 3.~The system returns review,
quiz and flashcard statistics and daily activity (\texttt{/learning/statistics/*},
\texttt{/learning/activity}), with days without activity filled with zeros. 4.~The interface
draws charts for 7, 30 or 90 days with Recharts.}
{3a.~The student views progress per document, per concept or per study session
(\texttt{/learning/progress/*}, \texttt{/learning/sessions}). 4a.~The student switches to the
table view.}
{E1.~No activity yet: the charts show zero values. E2.~Invalid or expired token: 401.}

\rclUseCase{10}{UC-10}{Manage AI providers (BYOK)}
{Primary: Student. Secondary: AI provider (key test)}
{The student is signed in.}
{The provider key is stored encrypted, its validity is recorded, and the chosen provider is
used for later AI requests.}
{1.~The student opens \emph{Settings}, tab \emph{AI \& Providers}. 2.~The student enters a key
(or an Ollama host URL) for Gemini, OpenAI, Groq or Ollama. 3.~\texttt{POST /api/v1/account/byok/credentials}
checks that the plan allows BYOK and that the provider is supported. 4.~The system encrypts the
key with AES-256-GCM and stores the ciphertext, nonce and tag with a hint of the last four
characters. 5.~The system tests the key against the provider's model list and records the
result. 6.~The student sets the provider as active.}
{2a.~The student tests a saved key without re-entering it (\texttt{POST /account/byok/test}).
2b.~The student deletes a key (\texttt{DELETE /account/byok/credentials/\{provider\}}).}
{E1.~Unsupported provider: 422. E2.~Plan does not allow BYOK: 403. E3.~Key rejected by the
provider: saved but marked invalid, with the provider's HTTP status shown.}

\rclUseCase{11}{UC-11}{Edit system prompts}
{Primary: Student (advanced user)}
{The application is running.}
{The custom template is used for every later request of that kind until it is reset.}
{1.~The student opens the LLM inspection page. 2.~\texttt{GET /api/prompts} lists the eleven
registered prompts (segment, flashcard, summary, explain, selection flashcard, chat, Socratic
question, Markdown clean-up, Cornell scaffold, Socratic evaluation, handwritten-outline
reconcile) with their variables and whether each is customised. 3.~The student edits a template.
4.~\texttt{PUT /api/prompts/\{key\}} stores it in the \texttt{system\_prompts} table. 5.~Later
calls to \texttt{get\_prompt\_template} return the custom text.}
{4a.~The student resets one prompt (\texttt{/api/prompts/\{key\}/reset}) or all prompts
(\texttt{/reset-all}). 1a.~The same page shows token usage (\texttt{/api/token-usage/summary},
\texttt{/history}) and can clear it.}
{E1.~Unknown prompt key, or a template whose placeholder syntax fails validation: 400, nothing
saved. E2.~At run time, a placeholder without a value is left in place instead of raising an
error (\texttt{SafePromptTemplate}).}

\rclUseCase{12}{UC-12}{Generate a flashcard set}
{Primary: Student. Secondary: AI provider}
{At least one processed document exists.}
{A named flashcard set with source links is saved and its cards enter the review queue.}
{1.~The student opens \emph{Flashcards} and chooses \emph{Generate}. 2.~The student selects
documents and a card count (1--50). 3.~\texttt{POST /api/v1/flashcards/generate} retrieves
chunks, reserves usage and asks the provider for cards. 4.~The system validates the cards and
saves the set with each card's source document and page. 5.~The student studies the set
(\texttt{/flashcards/:setId/study}); each card links back to its source page.}
{1a.~The student generates cards from a single topic (\texttt{POST /topics/\{id\}/flashcards}) or
from a highlighted passage (UC-03). 4a.~The student edits or deletes cards or the whole set.}
{E1.~Quota exceeded: 402. E2.~No usable chunks: 422. E3.~Provider error: generation fails and
the reserved usage is released.}

\clearpage
\section{Activity Diagram}

Four activity diagrams show the main workflows. \Cref{fig:act-study} gives the student's
overall study path, from signing in and checking that an AI provider is active, through import
and reading, to a choice between chat, drill, flashcard generation and quiz, and finally FSRS
review and progress analytics. It ties together UC-01 to UC-09.

\figplaceholder[0.8\textwidth]{activity_study.pdf}{Activity diagram of a student's study
workflow}{fig:act-study}

\Cref{fig:act-import} details UC-02. After the size and hash checks, the flow branches twice:
on whether a book with the same hash already exists for the user, and on whether the PDF is
scanned. A text PDF gets its topics from the built-in outline or from heading detection; a
scanned PDF goes to Marker if it is available (this OCR branch exists only on the book path). The
book is then mirrored into the application database, where it is chunked and embedded.

\figplaceholder[0.8\textwidth]{activity_import.pdf}{Activity diagram of importing and
processing a PDF (UC-02)}{fig:act-import}

\Cref{fig:act-review} details UC-06. The session loop fetches the next due card, shows the
question, waits for the reveal and the rating, and lets the FSRS scheduler compute the next due
date, until the queue is empty or the student abandons the session.

\figplaceholder[0.75\textwidth]{activity_review.pdf}{Activity diagram of reviewing due
flashcards with FSRS (UC-06)}{fig:act-review}

\Cref{fig:act-drill} details UC-05 and shows how a wrong explanation becomes study material:
the mastery score decides whether a nudge is given and how many flashcards are proposed, and
only the cards the student accepts are saved.

\figplaceholder[0.8\textwidth]{activity_drill.pdf}{Activity diagram of a Socratic drill,
from misconception to flashcards (UC-05)}{fig:act-drill}

\section{Static Model -- Class Diagram}

The static model is split into three class diagrams that follow the service classes of
\texttt{app/services}. Attributes are omitted where they add nothing; operations are the public
methods that the API routes call.

\Cref{fig:class-chat} shows how a chat answer is produced. \texttt{RAGService} uses
\texttt{RetrievalService} (composed of \texttt{VectorSearcher}, \texttt{KeywordSearcher} and a
\texttt{DiversityFilter}), \texttt{AIService} and \texttt{CitationValidator}, and asks
\texttt{EntitlementService} to reserve and finalise credits. \texttt{AIService} picks one
concrete \texttt{BaseAIAdapter}: \texttt{GeminiAdapter}, \texttt{OpenAIAdapter} (also used for Groq
and OpenRouter through their OpenAI-compatible APIs), \texttt{OllamaAdapter} or
\texttt{MockAIAdapter} for tests.

\widefig{class_chat_ai.pdf}{Class diagram: chat, retrieval and AI
providers}{fig:class-chat}

\Cref{fig:class-docs} shows document processing. \texttt{DocumentProcessingPipeline} reads the
file through \texttt{StorageService}, validates and extracts it with \texttt{ExtractionService}
(which returns an \texttt{ExtractedDocument} of pages), splits it with \texttt{SemanticChunker} and
embeds it with \texttt{EmbeddingService}. \texttt{BookSyncService} mirrors books from the study
database into the application database, and \texttt{KnowledgeService} builds document summaries
and concepts from the stored chunks.

\widefig{class_documents.pdf}{Class diagram: document upload and
processing}{fig:class-docs}

\Cref{fig:class-study} shows the study and billing services. \texttt{ReviewSessionService}
depends on the abstract \texttt{ReviewScheduler}, whose default implementation is
\texttt{FSRSScheduler}; \texttt{GraduatedIntervalScheduler} is used only if the \texttt{fsrs}
library is missing. \texttt{BillingService} depends on the abstract \texttt{BaseBillingProvider},
and \texttt{MockBillingProvider} implements it in the delivered build, so a payment provider
can later be added behind the same interface.

\widefig{class_study_billing.pdf}{Class diagram: study features and
billing}{fig:class-study}

\section{Dynamic Model -- Sequence Diagram}

\Cref{fig:seq-upload} shows \texttt{POST /api/v1/documents/upload}. The route checks the plan's
document and storage limits, saves the file, records it in both databases (reusing a book with
the same hash), creates the document and its processing job, and returns at once; the pipeline
then runs as a background task while the interface polls the job status.

\widefig{seq_upload.pdf}{Sequence diagram of uploading a document
(UC-02, \texttt{/api/v1} path)}{fig:seq-upload}

\Cref{fig:seq-chat} shows a streamed Knowledge Hub answer (UC-07): ownership check, provider
resolution, credit reservation, retrieval of five chunks, prompt assembly, token streaming as
SSE, citation validation and saving of the answer.

\widefig{seq_chat.pdf}{Sequence diagram of asking the Knowledge Hub
(UC-07)}{fig:seq-chat}

\Cref{fig:seq-drill} shows a Socratic drill (UC-05) on the study routes: question generation
with \texttt{socratic\_question\_prompt}, the student's answer, grading with
\texttt{socratic\_evaluation\_prompt} (registered in \texttt{app/prompt\_manager.py} and used by
\texttt{app/socratic\_drill.py}), and recording of the attempt and mastery score. These routes
do not pass through \texttt{EntitlementService}.

\widefig{seq_drill.pdf}{Sequence diagram of a Socratic drill
(UC-05)}{fig:seq-drill}

\Cref{fig:seq-login} shows UC-01: local login through \texttt{/api/v1/auth/login} (or Supabase
when configured), followed by \texttt{POST /api/set-user}, which switches the study database
file. The switch applies to the whole backend process, which is acceptable for a single-user
desktop application but would not be for a shared server.

\widefig{seq_login.pdf}{Sequence diagram of logging in
(UC-01)}{fig:seq-login}

\section{Safety and Security Requirements}\label{sec:security}

The requirements below are derived from \texttt{app/core/\allowbreak{}security.py}, \texttt{app/api/\allowbreak{}deps.py},
\texttt{app/api/\allowbreak{}middleware.py}, \texttt{app/core/\allowbreak{}config.py} and the BYOK routes in
\texttt{app/api/v1/\allowbreak{}account.py}, and follow the OWASP Top~10 guidance on broken access control,
cryptographic failures and authentication failures \cite{owasp2021}. Each requirement listed is
met by the delivered build.

\subsection{Access Requirements}

\begin{itemize}
  \item \textbf{SEC-A1 Authentication on the application API (Met).} Every \texttt{/api/v1}
  route except sign-up, login, the plan list, the billing webhook and the health checks requires
  a Bearer JWT, resolved by \texttt{get\_current\_user}. Tokens are signed with HS256, carry
  \texttt{sub}, \texttt{workspace\_id}, \texttt{iss=recall-ai}, \texttt{iat} and \texttt{exp},
  and expire after 60 minutes \cite{rfc7519}. Missing, expired or tampered tokens give 401.
  \item \textbf{SEC-A2 Workspace isolation (Met).} Data access is scoped to the caller's
  workspace (\texttt{get\_current\_workspace}). \texttt{require\_document\_owner} returns 404, not
  403, for a document in another workspace, so the existence of other users' records is not
  revealed (IDOR protection).
  \item \textbf{SEC-A3 Brute-force protection (Met).} Sign-up, login and change-password are
  limited to 15 requests per minute per client address by a sliding-window limiter; excess
  requests get 429. Login gives one message for an unknown e-mail and a wrong password.
  \item \textbf{SEC-A4 Disabled accounts (Met).} A user whose \texttt{is\_active} flag is false is
  rejected on every request, not only at login.
  \item \textbf{SEC-A5 Network exposure (Met).} The backend binds to \texttt{127.0.0.1} only (the
  sidecar's default host and \texttt{run.bat}), and CORS allows only the local development and
  Tauri origins.
\end{itemize}

\Cref{tab:access} summarises who can reach which \texttt{/api/v1} functions.

\begin{table}[htbp]
\centering
\caption{Access matrix of the \texttt{/api/v1} functions}\label{tab:access}
\small
\begin{tabularx}{\textwidth}{@{}X c c c@{}}
\toprule
\textbf{Function} & \makecell{\textbf{No token}} & \makecell{\textbf{Valid token,}\\\textbf{own data}} & \makecell{\textbf{Valid token,}\\\textbf{other user's data}}\\
\midrule
Sign up, log in, plan list, health & \cmark & \cmark & --\\
\texttt{/api/v1} documents, chat, flashcards, quizzes, reviews, progress & & \cmark & (404)\\
\texttt{/api/v1/account} and BYOK credentials & & \cmark & (404)\\
\bottomrule
\end{tabularx}
\end{table}

\subsection{Integrity Requirements}

\begin{itemize}
  \item \textbf{SEC-I1 Input validation (Met).} Every \texttt{/api/v1} request body is validated by
  Pydantic schemas (e-mail format, password length 8--128, counts 1--50, enumerated review
  ratings). Invalid input returns 422 with field details.
  \item \textbf{SEC-I2 Upload validation (Met).} Only PDF, Markdown and text files are accepted;
  a PDF must start with the \texttt{\%PDF} signature; PDFs are limited to 50~MB, uploads to
  55~MB and other request bodies to 10~MB (\texttt{RequestSizeLimitMiddleware}); oversized
  requests get 413. File names on disk are server-computed SHA-256 hashes, which prevents path
  traversal.
  \item \textbf{SEC-I3 Database consistency (Met).} Both SQLite databases run in WAL mode with
  \texttt{foreign\_keys=ON}; the application database also sets a 5{,}000~ms busy timeout. Sign-up
  creates user, workspace and preferences in one transaction.
  \item \textbf{SEC-I4 Server-side scoring (Met).} Quiz attempts are scored deterministically on
  the server, never by the client or the model.
  \item \textbf{SEC-I5 Validated AI output (Met).} Structured model output is repaired
  (\texttt{json-repair} and a truncation-repair pass) and validated against a schema before it is
  stored; citations are kept only if they point to retrieved chunks.
  \item \textbf{SEC-I6 Tamper-evident key storage (Met).} AES-GCM's 16-byte authentication tag is
  verified on every decryption, so a modified ciphertext is rejected \cite{nist80038d}.
  \item \textbf{SEC-I7 Webhook signatures (Met).} Billing webhooks are verified with HMAC-SHA256
  and a constant-time comparison; test signatures are accepted only when
  \texttt{ENVIRONMENT=test}.
\end{itemize}

\subsection{Privacy Requirements}

RecaLL processes personal data (name, e-mail, uploaded documents, notes, answers and study
history). It must follow Bangladesh's Personal Data Protection Ordinance 2025 (PDPO 2025) \cite{pdpo2025} and,
for any EU user, the principles of the GDPR \cite{gdpr2016}.

\begin{itemize}
  \item \textbf{SEC-P1 Local-first storage (Met).} All data is stored on the student's computer in
  the per-user application-data folder (\texttt{platformdirs}, application name ``Recall''); no
  RecaLL server receives it.
  \item \textbf{SEC-P2 Password storage (Met).} Passwords are stored only as
  \texttt{pbkdf2\_sha256\$600000\$salt\$hash} with a 16-byte random salt, matching the OWASP
  recommendation of 600{,}000 iterations for PBKDF2-HMAC-SHA256 \cite{owaspPassword}, and are
  checked with \texttt{hmac.compare\_digest}.
  \item \textbf{SEC-P3 Encrypted key vault (Met).} Keys saved through the v1 BYOK vault
  (\texttt{/api/v1/account/byok/credentials}, used by the provider cards under Settings $>$ AI
  Providers) are encrypted with AES-256-GCM and a fresh 12-byte nonce. The 32-byte master key is
  generated per installation and stored in the data folder (\texttt{.local\_byok\_key.hex}) unless
  set in the environment, and production mode refuses the default key. The vault only ever returns
  a hint (the last four characters).
  \item \textbf{SEC-P4 Data minimisation towards AI providers (Met).} Only the text a feature needs
  is sent: the selected chunks or topic text, the question or answer and, for chat, at most six
  earlier messages. No name, e-mail or account data is sent. With Ollama, nothing leaves the
  computer.
  \item \textbf{SEC-P5 Transparency (Met).} The settings page shows which provider is active, and
  the inspection page shows every prompt template and the recorded token usage. Langfuse tracing
  runs only if the student enters Langfuse keys.
  \item \textbf{SEC-P6 Secure defaults (Met).} Production mode refuses to start with the default
  JWT secret or BYOK key. Every response carries \texttt{X-Content-Type-Options: nosniff},
  \texttt{X-Frame-Options: DENY} and \texttt{Referrer-Policy:} \texttt{strict-origin-}\allowbreak\texttt{when-cross-origin}.
  \item \textbf{SEC-P7 Data deletion (Met).} The student can delete individual documents, notes,
  cards, quizzes and conversations, and can delete all study data and documents at once from
  Settings $>$ Security.
\end{itemize}

\section{AI-integration Requirements}\label{sec:ai-req}

RecaLL integrates hosted and local large language models instead of training its own, so its AI
requirements concern behaviour, grounding, robustness, latency, cost and human control, not
classification accuracy. \Cref{tab:airec} lists them. Values that come directly from the code
are fixed by design and verified by inspection or unit tests; values that need measurement are
marked and are reported in \Cref{ch:test}.

\begin{small}\setstretch{1.1}
\begin{longtable}{@{}l >{\raggedright\arraybackslash}p{7.4cm} c >{\raggedright\arraybackslash}p{3.6cm}@{}}
\caption{AI-integration requirements}\label{tab:airec}\\
\toprule
\textbf{ID} & \textbf{Requirement} & \textbf{Must/Want} & \textbf{Metric}\\
\midrule
\endfirsthead
\toprule
\textbf{ID} & \textbf{Requirement} & \textbf{Must/Want} & \textbf{Metric}\\
\midrule
\endhead
\bottomrule
\endfoot
R-AI-01 & The student shall be able to choose Gemini, OpenAI, Groq or a local Ollama model (and
OpenRouter through server settings) without changing code & Must & 4 providers selectable;
switch test\\
R-AI-02 & Requests made with the student's own key shall consume 0 RecaLL credits; the free plan
allows 50 credits per month & Must & Unit test on credit reservation\\
R-AI-03 & A drill shall contain exactly 2 questions: the first a causal mechanism, the second a
counterfactual or applied scenario; no recall-only questions & Must & Met: 3 of 3 topics gave valid
output with exactly 2 questions of the required types\\
R-AI-04 & Grading shall return a 0--100 mastery score in four bands (85--100, 60--84, 35--59,
0--34), with 0 cards at $\geq 85$, 1 card at 60--84 and 1--2 cards below 60, and no nudge below
60 & Must & Met: 4 of 4 graded answers followed the bands (strong answers 95 and 100 with 0 cards;
weak answers 0 and 0 with 2 cards each)\\
R-AI-05 & Every citation shown to the student shall point to a chunk retrieved for that answer
& Must & Fabricated citations shown = 0 (validator test)\\
R-AI-06 & Chat context shall be limited to 5 retrieved chunks, 3{,}000 tokens of reference text
and 6 earlier messages & Must & Code inspection / unit test\\
R-AI-07 & Document text shall be passed as isolated \texttt{<reference\_data>} blocks with closing
tags sanitised, and the model told not to follow instructions inside them \cite{owaspllm2025} & Must & Met:
delimiter-sanitisation test passes\\
R-AI-08 & Transient provider errors (HTTP 429, 502, 503, 504) shall be retried up to 3 times
with exponential back-off (2~s start, factor 2, cap 20~s, jitter up to 0.5~s); other HTTP errors,
such as 401 for a bad key, shall fail fast & Must & Unit test\\
R-AI-09 & Calls shall respect per-provider rate limits (Gemini free 5, Groq 30, Ollama 60,
OpenAI 500 requests per minute) & Must & Unit test\\
R-AI-10 & Structured outputs (drill questions, grading, cards, quizzes) shall be repaired and
schema-validated before saving; nothing invalid shall be stored & Must & Met: 0 parse failures in
7 structured replies (3 drills, 4 gradings)\\
R-AI-11 & A Knowledge Hub answer with a cloud provider shall be complete within 60~s & Want & Met:
4.9--44.8~s, median 15.9~s over 6 questions\\
R-AI-12 & A Cornell note for a typical topic shall be generated within 60~s & Want & Met:
26.2--47.9~s, median 27.9~s over 3 topics\\
R-AI-13 & Every prompt shall be editable and resettable by the student, without code changes &
Want & 11 of 11 prompts editable\\
R-AI-14 & AI output shall never be saved to the review deck without a student action (the student
selects which suggested cards to keep) & Must & Design review\\
R-AI-15 & Token usage of every call shall be recorded and visible to the student; tracing to
Langfuse shall be opt-in & Want & Usage visible on the inspection page\\
R-AI-16 & Retrieval should rank chunks by semantic similarity using model embeddings & Want &
Model embeddings (OpenAI, Gemini or Ollama) when a server key or Ollama host is configured\\
R-AI-17 & Scanned or handwritten PDFs shall be readable through OCR & Want & Optional Marker extractor
(local or Datalab API) on the book and topic path\\
\end{longtable}
\end{small}

\paragraph{Ethical requirements.} The AI only advises and drafts; the student remains in control
of what is kept, edited or revised (R-AI-14). The student can see and change every prompt
(R-AI-13), see which provider is used and choose a local model so that no material leaves the
computer (SEC-P4). The system must not present generated content as authoritative: cards and
notes stay linked to their source page so that claims can be checked. A full model card with the
models, prompts and usage notes is given in Appendix~F.
